\chapter{Flexible Cystoscopy (Male)}
\index{Flexible Cystoscopy (Male)}

\section*{Definition and Overview}
A flexible cystoscopy is a procedure in which a thin, bendable telescope is passed along the urethra into the bladder so that the lining of the bladder and urethra can be examined. In men, it is used to investigate blood in the urine, frequent or painful urination, recurrent urinary tract infections, difficulty passing urine, and bladder tumours, and to monitor the bladder after treatment for cancer. Because the male urethra is longer and narrower than the female, the procedure is usually slightly more uncomfortable, and local anaesthetic gel is used to numb the passage. It is a safe, quick outpatient procedure, usually lasting 5--10 minutes, and most men return to normal activities the same day. This chapter explains why it is done, what to expect, and how to recover.

\section*{Epidemiology}
Flexible cystoscopy is one of the most common urological procedures in many countries. It is performed frequently to investigate haematuria (blood in the urine), which affects around 1 in 10 adults, and to diagnose and survey bladder cancer, which is about three times more common in men than women and is strongly associated with smoking. Men also undergo cystoscopy for lower urinary tract symptoms from prostate problems, recurrent infection, and urethral conditions. Complications are uncommon, and the procedure is generally very well tolerated, though transient burning on urination is common afterwards.

\section*{Aetiology and Pathophysiology}
\begin{itemize}
    \item \textbf{Visualisation:} the flexible scope with a camera and light allows the urologist to inspect the whole bladder lining and the full length of the urethra
    \item \textbf{The male urethra:} longer and with curves and narrow points, which is why the procedure is more uncomfortable in men than in women
    \item \textbf{Anaesthetic gel:} a lubricating gel containing local anaesthetic is instilled into the urethra to numb it
    \item \textbf{Bladder filling:} sterile water or saline is instilled so the lining separates and can be fully examined
    \item \textbf{Diagnosis:} tumours, stones, inflammation, strictures, and an enlarged prostate can all be seen
    \item \textbf{Biopsy:} suspicious tissue can be sampled through the scope
    \item \textbf{Routes and flow:} the prostate surrounds the urethra, so prostate enlargement can compress the passage and be assessed indirectly
    \item \textbf{Comfort techniques:} slow passage, anaesthetic gel, and the flexible instrument keep discomfort to a minimum
\end{itemize}

\section*{Clinical Features}
\begin{itemize}
    \item The procedure takes 5--10 minutes in an outpatient clinic or day unit, without general anaesthesia
    \item You lie on your back, the penis is cleaned, and gel is inserted into the urethra before the scope is passed
    \item A sensation of needing to pass urine, pressure, and some discomfort as the scope passes
    \item Most men cope comfortably; the doctor can pause if needed
    \item Results may be discussed immediately, with biopsy results following
    \item Afterwards: burning or stinging on urination for 24--48 hours, and sometimes a little blood in the urine
    \item Normal activities can usually be resumed the same or next day
    \item The symptoms that prompted the test, such as blood in the urine or poor flow, guide the diagnosis
\end{itemize}

\section*{Differential Diagnosis}
Cystoscopy distinguishes the causes of urinary symptoms in men.
\begin{itemize}
    \item Bladder cancer and other bladder tumours
    \item Bladder stones
    \item Urinary tract infection and prostatitis
    \item Urethral stricture, from injury, infection, or previous procedures
    \item Benign prostatic enlargement, which narrows the urethra at the bladder neck
    \item Interstitial cystitis and other causes of bladder pain
    \item The source of blood in the urine, which may be in the kidneys and requires imaging as well
\end{itemize}

\section*{When to Seek Medical Care}
Men should see a doctor promptly for blood in the urine, persistent urinary symptoms, recurrent infections, or difficulty passing urine. Cystoscopy itself is an outpatient procedure, but these symptoms can signal serious conditions and need assessment without delay.

\textbf{Contact your local emergency services for:}
\begin{itemize}
    \item Inability to pass urine after the procedure (acute retention)
    \item Heavy bleeding, large clots, or prolonged bleeding
    \item Fever, chills, or severe pain, which may indicate infection or injury
    \item Severe lower abdominal pain
\end{itemize}

\section*{Diagnosis and Investigations}
\begin{itemize}
    \item \textbf{Urine tests:} dipstick, microscopy, and culture for blood, infection, and abnormal cells
    \item \textbf{Imaging:} ultrasound, CT urogram, or X-rays to examine the kidneys and ureters
    \item \textbf{The cystoscopy:} the definitive view of the bladder and urethra, with biopsy of suspicious areas
    \item \textbf{Histology:} pathology of any biopsy samples
    \item \textbf{Prostate assessment:} blood tests (PSA) and digital rectal examination where prostate disease is suspected
    \item \textbf{Surveillance:} repeat cystoscopy at intervals for bladder cancer, which commonly recurs
\end{itemize}

\section*{Management}
\begin{itemize}
    \item \textbf{Before:} empty the bladder; the doctor advises on medicines and blood thinners
    \item \textbf{During:} local anaesthetic gel and a gentle technique keep discomfort low
    \item \textbf{After:} drink extra fluids to dilute urine and ease burning
    \item \textbf{Pain relief:} paracetamol if needed for the first day or two
    \item \textbf{Antibiotics:} sometimes given for detected infection or after the procedure
    \item \textbf{Treating findings:} bladder tumours may be removed under anaesthetic; strictures may need dilatation; infection is treated with antibiotics
    \item \textbf{Follow-up:} discussion of results and a monitoring or treatment plan
\end{itemize}

\section*{Complications}
\begin{itemize}
    \item Urinary tract infection in a small proportion of men
    \item Blood in the urine for a day or two
    \item Burning on urination, usually mild and short-lived
    \item Difficulty passing urine, particularly in men with prostate enlargement
    \item Rarely, injury to the urethra or bladder
\end{itemize}

\section*{Prognosis and Outcomes}
Flexible cystoscopy in men is safe, quick, and generally well tolerated, and complications are uncommon. Its main value is diagnostic: it detects bladder tumours at an early, treatable stage, identifies stones and strictures, and gives reassurance when the bladder is normal. For men with bladder cancer, regular surveillance cystoscopy is essential because the disease recurs in a large proportion of patients. The short discomfort of the test is a small price for an accurate answer, and most men are back to normal the same day.

\section*{Prevention and Self-Care}
\begin{itemize}
    \item Understand the test and its purpose before consenting, and ask questions
    \item Mention all medicines, especially blood thinners
    \item Empty the bladder before the procedure
    \item Drink plenty of fluids afterwards
    \item Report fever, heavy bleeding, inability to pass urine, or severe pain
    \item Attend follow-up and surveillance appointments
    \item Do not ignore blood in the urine; seek assessment promptly
\end{itemize}

\section*{Sources and Further Reading}
The following reputable sources provide further information for healthcare students and patients seeking reliable, current advice about flexible cystoscopy in men.
\begin{itemize}
    \item Healthdirect Australia. \textit{Cystoscopy.} https://www.healthdirect.gov.au/cystoscopy
    \item The Urological Society of Australia and New Zealand. \textit{Cystoscopy patient information.} https://www.usanz.org.au/
    \item Better Health Channel. \textit{Cystoscopy.} https://www.betterhealth.vic.gov.au/
    \item NHS. \textit{Cystoscopy.} https://www.nhs.uk/conditions/cystoscopy/
    \item Mayo Clinic. \textit{Cystoscopy.} https://www.mayoclinic.org/
    \item MedlinePlus. \textit{Cystoscopy.} https://medlineplus.gov/
\end{itemize}
